\chapter{Files, catalog entries, and Algol programs}\label{files-catalog-entries-and-algol-programs}

Begin with an idle \texttt{muser} FP created by
\manualref{create-and-inspect-a-process}{Create and inspect a process}.
Skip the optional stop/remove exercise until after this tutorial.
If you already ended \texttt{muser}, request attention, select \texttt{s},
check \texttt{list} and create it again with
\texttt{new muser size 100000 temp disc 200 10 run}.
Use a fresh disposable session if the example file names already exist.

These examples use temporary entries and need the native image's
\texttt{fromhost} and \texttt{tohost}. In a second host terminal, change to
the emulator's working directory and use a text editor to create the ASCII/LF
file \texttt{exchange/input} shown below. Ensure that \texttt{exchange/output}
does not already exist: host exports deliberately refuse to overwrite files.

\section{Import, copy, rename, and export}\label{import-copy-rename-and-export}
\index{Files!import and export}
\index{Files!copying and renaming}
\index{lookup}
\index{copy}
\index{rename}

For the display and editor examples later in this manual, use exactly these
two lines as the contents of host \passthrough{\lstinline!exchange/input!},
including the final LF newline:

\begin{commandexample}{Host text file: exchange/input (not a command)}
Monitor manual sample.
Second line.
\end{commandexample}

Enter the following lines in FP, one at a time:

\begin{commandexample}{Guest FP in muser: import, copy, rename and export}
lookup fromhost tohost slang algol
msrc=fromhost input
lookup msrc
mcopy=copy msrc
rename mcopy.moved
lookup moved
output=tohost moved
\end{commandexample}

\passthrough{\lstinline!lookup!} displays catalog descriptors for the named entries. The first line locates the installed utilities; \passthrough{\lstinline!msrc!} and \passthrough{\lstinline!moved!} are temporary files in this user process. Their descriptors include allocation, device, date, and scope information. The simulator's offline disk listing is not a substitute for FP \passthrough{\lstinline!lookup!}: it cannot see changes that exist only in a running session's volatile overlay.

\passthrough{\lstinline!name=command ...!} directs the result to an RC catalog entry, except that the left-hand name of \passthrough{\lstinline!tohost!} is the host destination. \passthrough{\lstinline!copy!} copies the text into a new entry. The copy may use fewer disk segments than the source; equal text does not imply identical allocation. \passthrough{\lstinline!rename old.new!} changes the catalog name: the dot between the two names is required. The final line creates host \passthrough{\lstinline!exchange/output!}, containing the original text with LF line endings. Transfer byte counts can include RC padding; they need not equal the host file length.

Avoid naming entries after installed utilities. These are ordinary file operations, not disk initialization: do not use historical catalog creation or reset commands to prepare this session.

\begin{operationnote}{Guest names and host names are different}
\passthrough{\lstinline!msrc!} and \passthrough{\lstinline!moved!} are
guest catalog names. \passthrough{\lstinline!input!} and
\passthrough{\lstinline!output!} name host files beneath the configured
exchange root. Do not enter \passthrough{\lstinline!exchange/input!} as a
guest catalog name or try to use a guest name as a host pathname. If a host
export already exists, choose a new output name instead of deleting unrelated
files or assuming the export replaced it.
\end{operationnote}

\section{Display a text file on the guest console}\label{display-a-text-file-on-the-guest-console}

In the same FP process, use the special console destination \passthrough{\lstinline!c!}:

\begin{commandexample}{Guest FP in muser: display one text file}
c=copy msrc
\end{commandexample}

For an imported ASCII/LF file containing \passthrough{\lstinline!Monitor manual sample.!} and \passthrough{\lstinline!Second line.!}, the two lines appear on the guest console. A \passthrough{\lstinline!segm.!} copy summary follows the text; it is a utility report, not part of the file. The command returns to FP without attention or a reset. It does not edit the source or create a new disk file named \passthrough{\lstinline!c!}.

To display several text files consecutively, name them in the desired order:

\begin{commandexample}{Guest FP in muser: display two text files}
c=copy msrc moved
\end{commandexample}

Since \passthrough{\lstinline!moved!} is the unchanged copy in the earlier example, this displays the same two-line text twice, followed by input and destination summaries. The tested source and copy remain unchanged after both displays.

The destination matters: \passthrough{\lstinline!result=copy source!} creates a disk result, whereas \passthrough{\lstinline!c=copy source!} sends text to the console. Do not assume bare \passthrough{\lstinline!copy source!} is a file-viewing command; it produced only a copy report in the exploratory session. Do not substitute host \passthrough{\lstinline!cat!}: the installed guest \passthrough{\lstinline!cat!} is an entry point in \passthrough{\lstinline!catsort!}, not the host text utility.

This is a text display, not a pager, binary dump or full-screen editor. Large files can scroll past the terminal's history; export them with \passthrough{\lstinline!tohost!} and inspect the host file when necessary. Additional \passthrough{\lstinline!copy!} options, other console destinations and control-character handling remain unverified by this example.

\section{Compile and run a small Algol program}\label{compile-and-run-a-small-algol-program}
\index{Algol}

The optional tutorial bundle includes this example as \passthrough{\lstinline!examples/hello.alg!}:

\begin{commandexample}{Algol source: hello.alg (not an FP command)}
begin
   write(out, <:Hello from Algol<10>:>)
end
\end{commandexample}

In a second host terminal, change to the same working directory as the
emulator and copy the bundled source. If you did not copy \texttt{examples/},
create \texttt{exchange/alginput} with exactly the source shown above instead.

\begin{commandexample}[language=sh]{Host shell in the emulator's working directory: prepare Algol input}
cp examples/hello.alg exchange/alginput
\end{commandexample}

Then enter these FP commands in the existing \passthrough{\lstinline!muser!} process:

\begin{commandexample}{Guest FP in muser: compile and run Algol}
asrc=fromhost alginput
mprog=algol asrc connect.no
mprog
\end{commandexample}

The supplied compiler prints \passthrough{\lstinline!algol end 24!} for this source. Running \passthrough{\lstinline!mprog!} prints \passthrough{\lstinline!Hello from Algol!}, then returns to FP with \passthrough{\lstinline!end 7!}. The end numbers are program reports, not host exit statuses. Compilation and execution both happen on the emulated RC4000, not through the host C compiler.

\section{Clean up and keep results}\label{clean-up-and-keep-results}
\index{clear}

If you want to continue with the
\manualref{edit-a-text-file-inside-the-guest}{editor} or
\manualref{compare-text-files-inside-the-guest}{text comparison} examples,
keep \passthrough{\lstinline!msrc!} and \passthrough{\lstinline!moved!}
until those examples are complete. Skip cleanup for now if you will continue
with those recipes; their prerequisites also explain how to recreate the files.
The cleanup below deliberately removes both names.

After exporting any wanted results, remove just the example entries in FP:

\begin{commandexample}{Guest FP in muser: remove the example entries}
clear temp msrc moved asrc mprog
\end{commandexample}

This deletes those temporary catalog entries. Successful \passthrough{\lstinline!rename!} and \passthrough{\lstinline!clear!} need not print a confirmation. Use \passthrough{\lstinline!lookup!} to inspect entries when unsure, and check the current command context before retrying a command that reports a syntax error or cannot find a program.

Temporary catalog scope and emulator disk persistence are separate concepts. A persistent overlay or writable image does not turn a temporary entry into a permanent file: the guest can reclaim it when its process is removed. For the workflow above, \passthrough{\lstinline!tohost!} is the verified way to keep text beyond the guest session. See \manualref{keep-guest-files-across-sessions}{the persistent guest-file workflow}; advanced catalog administration and the full historical utility set are not covered by these examples. Bounded command-file examples are described separately in \manualref{call-one-command-file-from-another}{FP command-file input}.
